% rel_framework_mhd.tex
% Relativistic MHD Framework
% Conventions: (-,+,+,+) signature, c kept explicit, u^μ u_μ = -c²
% See RELATIVISTIC_CONVENTIONS.md for notation and macros.

\chapter{Relativistic Magnetohydrodynamic Framework}
\label{ch:rel-mhd-framework}

The classical theory of magnetohydrodynamic stability, as developed by
Chandrasekhar in Chapters~IV--VI, rests on the Navier--Stokes and
induction equations in the non-relativistic regime.  When the Alfv\'en
speed $\vA$ becomes an appreciable fraction of $c$, or when the fluid is
embedded in a strongly curved spacetime, a fully covariant formulation is
required.  This chapter develops the relativistic MHD framework from first
principles, keeping the speed of light $c$ explicit throughout so that
every non-relativistic limit can be read off transparently.

%======================================================================
\section{Covariant Electrodynamics}
\label{sec:covariant-ed}
%======================================================================

\subsection{The Faraday tensor and its dual}

The electromagnetic field is encoded in the antisymmetric Faraday tensor
$\farad$.  In a coordinate basis with the metric signature $(-,+,+,+)$,
its components read
\begin{equation}
  F^{\mu\nu} =
  \begin{pmatrix}
    0      & -E^{1}/c & -E^{2}/c & -E^{3}/c \\
    E^{1}/c &  0       & -B^{3}   &  B^{2}   \\
    E^{2}/c &  B^{3}   &  0       & -B^{1}   \\
    E^{3}/c & -B^{2}   &  B^{1}   &  0
  \end{pmatrix},
  \label{eq:Fmunu-components}
\end{equation}
where $E^{i}$ and $B^{i}$ are the electric and magnetic field components
measured by an observer at rest in the coordinate frame.

The dual tensor is defined by
\begin{equation}
  {}^{\ast}\!F^{\mu\nu}
  = \frac{1}{2}\,\epsilon^{\mu\nu\alpha\beta}\,F_{\alpha\beta},
  \label{eq:dual-Faraday}
\end{equation}
with $\epsilon^{\mu\nu\alpha\beta}$ the Levi-Civita tensor satisfying
$\epsilon^{0123} = 1/\!\sqrt{-g}$.

\subsection{Maxwell equations in covariant form}

The inhomogeneous Maxwell equations are
\begin{equation}
  \nabla_{\mu}\,F^{\mu\nu} = \frac{1}{c}\,j^{\nu},
  \label{eq:Maxwell-inhom}
\end{equation}
where $j^{\nu}$ is the charge--current 4-vector.  The homogeneous
equations (the Bianchi identity) take the form
\begin{equation}
  \nabla_{[\alpha}\,F_{\beta\gamma]} = 0
  \qquad\Longleftrightarrow\qquad
  \nabla_{\mu}\,{}^{\ast}\!F^{\mu\nu} = 0.
  \label{eq:Maxwell-hom}
\end{equation}
These two sets of equations, \eqref{eq:Maxwell-inhom} and
\eqref{eq:Maxwell-hom}, are manifestly covariant and hold in arbitrary
curved spacetimes without modification.

\subsection{Electric and magnetic fields in the fluid frame}

Given a fluid with 4-velocity $\fvel$ (normalised so that
$u^{\mu}u_{\mu}=-c^{2}$), one decomposes the electromagnetic field into
electric and magnetic parts as measured in the instantaneous rest frame:
\begin{equation}
  E^{\mu} = F^{\mu\nu}\,u_{\nu},
  \label{eq:Emu-def}
\end{equation}
\begin{equation}
  b^{\mu} = \frac{1}{c}\,{}^{\ast}\!F^{\mu\nu}\,u_{\nu}
  = \frac{1}{2c}\,\epsilon^{\mu\nu\alpha\beta}\,u_{\nu}\,F_{\alpha\beta}.
  \label{eq:bmu-def}
\end{equation}
Both $E^{\mu}$ and $b^{\mu}$ are orthogonal to $u^{\mu}$:
\begin{equation}
  E^{\mu}\,u_{\mu} = 0, \qquad b^{\mu}\,u_{\mu} = 0.
  \label{eq:E-b-orthogonal}
\end{equation}
The scalar $b^{2}\equiv b^{\mu}b_{\mu}$ plays the role of twice the
magnetic pressure, and equals $B^{2}$ (the square of the rest-frame
magnetic field magnitude) up to factors of $4\pi$.

The Faraday tensor can be reconstructed from these fluid-frame fields:
\begin{equation}
  F^{\mu\nu}
  = \frac{1}{c}\bigl(E^{\mu}\,u^{\nu} - E^{\nu}\,u^{\mu}\bigr)
  + \epsilon^{\mu\nu\alpha\beta}\,u_{\alpha}\,b_{\beta}/c.
  \label{eq:Fmunu-decomp}
\end{equation}


%======================================================================
\section{Ideal Relativistic MHD}
\label{sec:ideal-rmhd}
%======================================================================

\subsection{The ideal MHD condition}

In a perfectly conducting fluid, the electric field vanishes in the
comoving frame.  This \emph{ideal MHD condition} is
\begin{equation}
  E^{\mu} = F^{\mu\nu}\,u_{\nu} = 0.
  \label{eq:ideal-MHD-condition}
\end{equation}
Equivalently, the Faraday tensor reduces to
\begin{equation}
  F^{\mu\nu} = \frac{1}{c}\,\epsilon^{\mu\nu\alpha\beta}\,u_{\alpha}\,b_{\beta}.
  \label{eq:F-ideal}
\end{equation}

\subsection{Energy-momentum tensor of the magnetised fluid}

The total energy-momentum tensor is the sum of fluid and electromagnetic
contributions:
\begin{equation}
  T^{\mu\nu} = T^{\mu\nu}_{\mathrm{fluid}} + T^{\mu\nu}_{\mathrm{EM}}.
  \label{eq:Tmunu-total}
\end{equation}
For a perfect fluid,
\begin{equation}
  T^{\mu\nu}_{\mathrm{fluid}}
  = \frac{\varepsilon + p}{c^{2}}\,u^{\mu}\,u^{\nu} + p\,g^{\mu\nu},
  \label{eq:Tmunu-fluid}
\end{equation}
where $\varepsilon$ is the total energy density (including rest-mass
energy) and $p$ the isotropic pressure.

Under the ideal MHD condition \eqref{eq:ideal-MHD-condition}, the
electromagnetic energy-momentum tensor becomes
\begin{equation}
  T^{\mu\nu}_{\mathrm{EM}}
  = \frac{b^{2}}{c^{2}}\,u^{\mu}\,u^{\nu}
  + \frac{b^{2}}{2}\,g^{\mu\nu}
  - b^{\mu}\,b^{\nu}.
  \label{eq:Tmunu-EM}
\end{equation}
The first term represents the magnetic energy density's contribution to
inertia; the second is the isotropic magnetic pressure $b^{2}/2$; the
third is the anisotropic magnetic tension.

Combining \eqref{eq:Tmunu-fluid} and \eqref{eq:Tmunu-EM}, the total
energy-momentum tensor of the ideal relativistic MHD fluid is
\begin{equation}
  T^{\mu\nu}
  = \frac{\varepsilon + p + b^{2}}{c^{2}}\,u^{\mu}\,u^{\nu}
  + \Bigl(p + \frac{b^{2}}{2}\Bigr)\,g^{\mu\nu}
  - b^{\mu}\,b^{\nu}.
  \label{eq:Tmunu-MHD}
\end{equation}
This has the form of a perfect fluid with effective enthalpy density
$\enthalpy^{\ast} = (\varepsilon + p + b^{2})/c^{2}$ and effective
pressure $p^{\ast} = p + b^{2}/2$, modified by the anisotropic tension
$b^{\mu}b^{\nu}$.

\subsection{Covariant induction equation}

The evolution of the magnetic field is governed by the induction equation.
Taking the ideal MHD condition and the homogeneous Maxwell equation
\eqref{eq:Maxwell-hom}, one obtains
\begin{equation}
  u^{\mu}\,\nabla_{\mu}\,b^{\nu}
  - b^{\mu}\,\nabla_{\mu}\,u^{\nu}
  + b^{\nu}\,\nabla_{\mu}\,u^{\mu}
  - u^{\nu}\,\nabla_{\mu}\,b^{\mu}
  = 0.
  \label{eq:induction-covariant}
\end{equation}
Using the constraint $b^{\mu}u_{\mu}=0$ and the normalisation
$u^{\mu}u_{\mu}=-c^{2}$, this can be written more compactly as
\begin{equation}
  \mathcal{L}_{u}\,b^{\mu}
  = b^{\mu}\,\nabla_{\nu}\,u^{\nu}
    - b^{\nu}\,\nabla_{\nu}\,u^{\mu}
    + u^{\mu}\,b^{\nu}\,\frac{a_{\nu}}{c^{2}},
  \label{eq:induction-Lie}
\end{equation}
where $\mathcal{L}_{u}$ denotes the Lie derivative along $u^{\mu}$ and
$a^{\mu} = u^{\nu}\nabla_{\nu}u^{\mu}$ is the 4-acceleration.  In the
non-relativistic limit ($c\to\infty$) this reduces to the familiar form
$\partial\vect{B}/\partial t = \nabla\times(\vect{v}\times\vect{B})$.

\subsection{Magnetic flux freezing: the relativistic Alfv\'en theorem}

Alfv\'en's theorem---that magnetic flux through any surface comoving
with the fluid is conserved---generalises directly to the relativistic
regime.  Consider a 2-surface $\mathcal{S}$ transported by the fluid
flow.  The magnetic flux through $\mathcal{S}$ is
\begin{equation}
  \Phi_{\mathcal{S}}
  = \int_{\mathcal{S}} F_{\mu\nu}\,\dd\Sigma^{\mu\nu},
  \label{eq:flux-def}
\end{equation}
where $\dd\Sigma^{\mu\nu}$ is the surface element.  By virtue of the
ideal MHD condition \eqref{eq:ideal-MHD-condition} and the homogeneous
Maxwell equation \eqref{eq:Maxwell-hom}, one can show that
\begin{equation}
  \frac{\dd\Phi_{\mathcal{S}}}{\dd\tau} = 0,
  \label{eq:flux-freezing}
\end{equation}
where $\tau$ is the proper time along the fluid worldlines.  This is the
relativistic flux-freezing theorem: magnetic field lines are ``frozen
into'' the perfectly conducting fluid, exactly as in the non-relativistic
case.


%======================================================================
\section{Characteristic Speeds and Causality}
\label{sec:char-speeds}
%======================================================================

A central requirement of relativistic physics is that all signal speeds
remain bounded by $c$.  We now derive the characteristic speeds of ideal
relativistic MHD and verify that this is automatically satisfied.

\subsection{Alfv\'en speed}

The relativistic Alfv\'en speed is obtained from the eigenvalue analysis
of the linearised MHD system.  For propagation along the magnetic field
direction,
\begin{equation}
  \vA^{2}
  = \frac{b^{2}}{4\pi\enthalpy^{\ast}}
  = \frac{b^{2}}{4\pi\bigl[(\varepsilon + p)/c^{2} + b^{2}/c^{2}\bigr]}
  = \frac{b^{2}\,c^{2}}{4\pi(\varepsilon + p) + b^{2}},
  \label{eq:vA-rel}
\end{equation}
where we have restored factors of $4\pi$ for Gaussian units.  Since the
denominator exceeds the numerator whenever $\varepsilon + p > 0$ (as
guaranteed by the dominant energy condition), we have the fundamental bound
\begin{equation}
  \vA^{2} < c^{2}.
  \label{eq:vA-bound}
\end{equation}
In the non-relativistic limit $b^{2} \ll 4\pi\rho c^{2}$, this reduces to
the classical expression $\vA^{2} \to B^{2}/(4\pi\rho)$.

\subsection{Sound speed}

The adiabatic sound speed is
\begin{equation}
  \cs^{2} = \left(\frac{\partial p}{\partial\varepsilon}\right)_{\!s}\,c^{2},
  \label{eq:cs-rel}
\end{equation}
and thermodynamic stability (plus causality) requires
\begin{equation}
  0 \le \cs^{2} \le c^{2}.
  \label{eq:cs-bound}
\end{equation}

\subsection{Fast and slow magnetosonic speeds}

The fast and slow magnetosonic speeds for propagation at angle $\theta$
to the magnetic field are the roots of
\begin{equation}
  v^{4} - v^{2}\bigl(\cs^{2} + \vA^{2}
    - \cs^{2}\vA^{2}/c^{2}\bigr)
  + \cs^{2}\,\vA^{2}\,\cos^{2}\!\theta = 0.
  \label{eq:magnetosonic-dispersion}
\end{equation}
The solutions are
\begin{align}
  v_{\mathrm{f}}^{2}
  &= \frac{1}{2}\Bigl[\bigl(\cs^{2}+\vA^{2}
     -\cs^{2}\vA^{2}/c^{2}\bigr)
    + \sqrt{\Delta}\;\Bigr],
  \label{eq:vf}\\[6pt]
  v_{\mathrm{s}}^{2}
  &= \frac{1}{2}\Bigl[\bigl(\cs^{2}+\vA^{2}
     -\cs^{2}\vA^{2}/c^{2}\bigr)
    - \sqrt{\Delta}\;\Bigr],
  \label{eq:vs}
\end{align}
where
\begin{equation}
  \Delta
  = \bigl(\cs^{2}+\vA^{2}-\cs^{2}\vA^{2}/c^{2}\bigr)^{2}
    - 4\,\cs^{2}\,\vA^{2}\,\cos^{2}\!\theta.
  \label{eq:Delta}
\end{equation}

\subsection{Causality verification}

The fast magnetosonic speed is maximised for $\theta=\pi/2$ (propagation
perpendicular to $\vect{B}$), where $v_{\mathrm{s}}^{2}=0$ and
\begin{equation}
  v_{\mathrm{f}}^{2}\big|_{\theta=\pi/2}
  = \cs^{2} + \vA^{2} - \frac{\cs^{2}\,\vA^{2}}{c^{2}}.
  \label{eq:vf-perp}
\end{equation}
Since $\cs^{2}<c^{2}$ and $\vA^{2}<c^{2}$, we can write
\begin{equation}
  v_{\mathrm{f}}^{2}
  = \cs^{2} + \vA^{2}\!\left(1-\frac{\cs^{2}}{c^{2}}\right)
  < \cs^{2} + c^{2}\!\left(1-\frac{\cs^{2}}{c^{2}}\right)
  = c^{2}.
  \label{eq:vf-causal-proof}
\end{equation}
Thus $v_{\mathrm{f}}^{2} < c^{2}$ identically.  Since
$v_{\mathrm{s}}^{2}\le v_{\mathrm{f}}^{2}$ and
$\vA^{2}<c^{2}$, \emph{all three characteristic speeds of ideal
relativistic MHD are strictly subluminal}.  This is a fundamental
self-consistency of the relativistic theory that has no counterpart in the
non-relativistic framework, where $\vA$ is unbounded.

The ordering of characteristic speeds is
\begin{equation}
  0 \;\le\; v_{\mathrm{s}} \;\le\; \min(\cs,\,\vA)
  \;\le\; \max(\cs,\,\vA) \;\le\; v_{\mathrm{f}} \;<\; c.
  \label{eq:speed-ordering}
\end{equation}


%======================================================================
\section{Resistive Relativistic MHD}
\label{sec:resistive-rmhd}
%======================================================================

\subsection{Relativistic Ohm's law}

When the conductivity $\sigma$ is finite, the ideal MHD condition
\eqref{eq:ideal-MHD-condition} is replaced by the relativistic Ohm's law:
\begin{equation}
  j^{\mu} + \frac{j^{\nu}\,u_{\nu}}{c^{2}}\,u^{\mu}
  = \sigma\,F^{\mu\nu}\,u_{\nu}
  = \sigma\,E^{\mu}.
  \label{eq:rel-Ohm}
\end{equation}
The left-hand side projects the current density onto the rest frame of the
fluid (removing the charge density component $j^{\nu}u_{\nu}/c^{2}$),
and the right-hand side is the rest-frame electric field multiplied by the
scalar conductivity $\sigma$.  In the limit $\sigma\to\infty$, one
recovers $E^{\mu}\to 0$, i.e.\ the ideal MHD condition.

Written in terms of the projection tensor
$\proj = g^{\mu\nu}+u^{\mu}u^{\nu}/c^{2}$, this is
\begin{equation}
  \Delta^{\mu}{}_{\nu}\,j^{\nu} = \sigma\,E^{\mu}.
  \label{eq:rel-Ohm-proj}
\end{equation}

\subsection{Magnetic diffusivity}

The magnetic diffusivity is defined (in Gaussian units) by
\begin{equation}
  \eta_{\mathrm{m}} = \frac{c^{2}}{4\pi\sigma},
  \label{eq:eta-m}
\end{equation}
with dimensions of length$^{2}$/time.  This is the direct relativistic
generalisation of the classical magnetic diffusivity.

\subsection{Relativistic magnetic Reynolds number}

The magnetic Reynolds number, which measures the ratio of advective to
diffusive transport of the magnetic field, is
\begin{equation}
  \mathrm{Rm} = \frac{v\,L}{\eta_{\mathrm{m}}}
  = \frac{4\pi\sigma\,v\,L}{c^{2}},
  \label{eq:Rm-rel}
\end{equation}
where $v$ is a characteristic velocity and $L$ a characteristic length
scale.  The factor of $c^{2}$ in the denominator reflects the relativistic
bound: as $v\to c$, the maximum magnetic Reynolds number for given $\sigma$
and $L$ is $\mathrm{Rm}_{\max} = 4\pi\sigma L/c$.

For $\mathrm{Rm}\gg 1$, the field is effectively frozen into the fluid
and ideal MHD applies.  For $\mathrm{Rm}\lesssim 1$, resistive effects
dominate and field diffusion must be accounted for.

\subsection{Magnetic field decay and Joule dissipation}

In a resistive fluid, the induction equation acquires a diffusive term.
The covariant induction equation becomes
\begin{equation}
  u^{\mu}\,\nabla_{\mu}\,b^{\nu}
  - b^{\mu}\,\nabla_{\mu}\,u^{\nu}
  + b^{\nu}\,\nabla_{\mu}\,u^{\mu}
  = \eta_{\mathrm{m}}\,\Delta^{\nu\alpha}\,
    \nabla_{\beta}\,\nabla^{\beta}\,b_{\alpha}
  + \text{(curvature terms)},
  \label{eq:induction-resistive}
\end{equation}
where the right-hand side represents the magnetic diffusion (the
Laplacian term) projected onto the fluid rest frame.  In flat spacetime,
this reduces to
\begin{equation}
  \frac{\partial\vect{B}}{\partial t}
  = \nabla\times(\vect{v}\times\vect{B})
  + \eta_{\mathrm{m}}\,\nabla^{2}\vect{B},
  \label{eq:induction-NR}
\end{equation}
which is the standard non-relativistic induction equation with Ohmic
diffusion.

The Joule heating rate per unit volume in the fluid frame is
\begin{equation}
  Q_{\mathrm{Joule}}
  = \frac{E^{\mu}\,E_{\mu}}{\eta_{\mathrm{m}}}
  = \frac{j^{\mu}j_{\mu} + (j^{\nu}u_{\nu})^{2}/c^{2}}{\sigma},
  \label{eq:Joule}
\end{equation}
which enters as a source term in the energy equation.  In the ideal limit
$\sigma\to\infty$, both $E^{\mu}\to 0$ and $Q_{\mathrm{Joule}}\to 0$.

The characteristic timescale for Ohmic decay of a magnetic field
structure of size $L$ is
\begin{equation}
  \tau_{\mathrm{Ohm}} = \frac{L^{2}}{\eta_{\mathrm{m}}}
  = \frac{4\pi\sigma\,L^{2}}{c^{2}}.
  \label{eq:tau-Ohm}
\end{equation}


%======================================================================
\section{Linearised Relativistic MHD Perturbations}
\label{sec:linearised-rmhd}
%======================================================================

\subsection{Background equilibrium}

Consider a static, uniform background:
\begin{itemize}
  \item Fluid at rest: $u^{\mu}_{(0)} = (c,\,0,\,0,\,0)$ in Minkowski
        spacetime.
  \item Uniform energy density $\varepsilon_{0}$ and pressure $p_{0}$.
  \item Uniform magnetic field $\vect{B}_{0} = B_{0}\,\hat{\vect{z}}$,
        so that $b^{\mu}_{(0)} = (0,\,0,\,0,\,B_{0})$ in the rest frame,
        and $b^{2}_{(0)} = B_{0}^{2}$.
\end{itemize}

\subsection{Perturbation ansatz}

We perturb all quantities:
\begin{equation}
  u^{\mu} = u^{\mu}_{(0)} + \delta u^{\mu},
  \qquad
  \varepsilon = \varepsilon_{0} + \delta\varepsilon,
  \qquad
  p = p_{0} + \delta p,
  \qquad
  b^{\mu} = b^{\mu}_{(0)} + \delta b^{\mu},
  \label{eq:perturbation-ansatz}
\end{equation}
and seek plane-wave solutions $\sim \exp[i(k_{\mu}x^{\mu})]$ with
4-wavevector $k^{\mu} = (\omega/c,\,\vect{k})$.  The propagation angle
relative to the background field is $\theta$, so
$\vect{k}\cdot\hat{\vect{z}} = k\cos\theta$.

The normalisation constraint $u^{\mu}u_{\mu}=-c^{2}$ implies
$u_{(0)\mu}\,\delta u^{\mu} = 0$, so $\delta u^{0}=0$ at first order---the
velocity perturbation is purely spatial: $\delta u^{\mu} = (0,\,\delta\vect{v})$.

\subsection{Linearised equations}

Linearising the energy-momentum conservation $\nabla_{\mu}T^{\mu\nu}=0$
and the induction equation \eqref{eq:induction-covariant} about the
background, one obtains a system of linear algebraic equations (in
Fourier space):
\begin{align}
  -\omega\,\enthalpy^{\ast}_{0}\,\delta v^{i}
  &= k^{i}\,c^{2}\Bigl(\delta p + \frac{b_{(0)}^{j}\,\delta b_{j}}{4\pi}\Bigr)
  - \frac{c^{2}}{4\pi}\,(k_{j}\,b_{(0)}^{j})\,\delta b^{i}
  \notag\\
  &\quad
  - \frac{c^{2}}{4\pi}\,(k_{j}\,\delta b^{j})\,b_{(0)}^{i},
  \label{eq:lin-momentum}
\end{align}
\begin{equation}
  -\omega\,\delta\varepsilon
  = (\varepsilon_{0}+p_{0})\,(k_{j}\,\delta v^{j}),
  \label{eq:lin-energy}
\end{equation}
\begin{equation}
  -\omega\,\delta b^{i}
  = (k_{j}\,b_{(0)}^{j})\,\delta v^{i}
  - b_{(0)}^{i}\,(k_{j}\,\delta v^{j}),
  \label{eq:lin-induction}
\end{equation}
where $\enthalpy^{\ast}_{0} = (\varepsilon_{0}+p_{0}+b_{(0)}^{2})/(4\pi c^{2})$
is the background total (relativistic) inertia, with appropriate factors
of $4\pi$ restored for Gaussian units.

\subsection{Dispersion relation}

Eliminating $\delta\varepsilon$, $\delta p = \cs^{2}\,\delta\varepsilon/c^{2}$,
and $\delta b^{i}$ in favour of $\delta v^{i}$, one arrives at the
dispersion relation.  Defining the phase speed $v_{\mathrm{ph}}=\omega/k$,
the result is:
\begin{equation}
  \bigl(v_{\mathrm{ph}}^{2} - \vA^{2}\cos^{2}\!\theta\bigr)
  \Bigl[v_{\mathrm{ph}}^{4}
    - v_{\mathrm{ph}}^{2}\bigl(\cs^{2}+\vA^{2}
      -\cs^{2}\vA^{2}/c^{2}\bigr)
    + \cs^{2}\,\vA^{2}\,\cos^{2}\!\theta\Bigr] = 0.
  \label{eq:dispersion-full}
\end{equation}
This factors into:
\begin{enumerate}
  \item \textbf{Alfv\'en mode}: $v_{\mathrm{ph}}^{2} = \vA^{2}\cos^{2}\!\theta$,
  \item \textbf{Fast/slow magnetosonic modes}: the quartic in
        $v_{\mathrm{ph}}^{2}$ whose roots are given by
        \eqref{eq:vf}--\eqref{eq:vs}.
\end{enumerate}
All three families of modes satisfy $v_{\mathrm{ph}}^{2} < c^{2}$, as
proved in \S\ref{sec:char-speeds}.

\subsection{Non-relativistic limit: recovery of Chandrasekhar Ch~IV}

In the non-relativistic limit, $\varepsilon \to \rho c^{2}$,
$p \ll \rho c^{2}$, $b^{2}\ll 4\pi\rho c^{2}$, and $\cs^{2}\ll c^{2}$.
The relativistic dispersion \eqref{eq:dispersion-full} then reduces to
\begin{equation}
  \bigl(v_{\mathrm{ph}}^{2} - V_{\!A}^{2}\cos^{2}\!\theta\bigr)
  \bigl[v_{\mathrm{ph}}^{4}
    - v_{\mathrm{ph}}^{2}(a^{2}+V_{\!A}^{2})
    + a^{2}\,V_{\!A}^{2}\,\cos^{2}\!\theta\bigr] = 0,
  \label{eq:dispersion-NR}
\end{equation}
where $V_{\!A}^{2} = B^{2}/(4\pi\rho)$ and $a^{2}=(\partial p/\partial\rho)_{s}$
are the non-relativistic Alfv\'en and sound speeds.  This is precisely the
dispersion relation derived in Chandrasekhar~\cite{Chandrasekhar1961},
Chapter~IV, equations~(38) and~(39).


%======================================================================
\section{Relativistic Alfv\'en Waves}
\label{sec:rel-alfven}
%======================================================================

\subsection{Derivation of the Alfv\'en mode}

The Alfv\'en mode corresponds to incompressible, transverse perturbations
with $\delta\varepsilon=0$, $\vect{k}\cdot\delta\vect{v}=0$.  From the
linearised equations \eqref{eq:lin-momentum}--\eqref{eq:lin-induction},
the perturbation velocity is perpendicular to both $\vect{k}$ and
$\vect{B}_{0}$ (for general $\theta$), and the dispersion relation is
simply
\begin{equation}
  \omega^{2} = k^{2}\,\vA^{2}\,\cos^{2}\!\theta,
  \label{eq:alfven-dispersion}
\end{equation}
with the relativistic Alfv\'en speed given by \eqref{eq:vA-rel}.  The
corresponding phase velocity along the magnetic field is
\begin{equation}
  v_{\mathrm{ph}} = \vA\cos\theta
  = \cos\theta\;\frac{b\,c}{\sqrt{4\pi(\varepsilon+p)+b^{2}}},
  \label{eq:alfven-phase}
\end{equation}
where $b = |\vect{B}_{0}|$ in the rest frame.  This is manifestly less
than $c$ for all $\theta$.

The group velocity for Alfv\'en waves is directed along the background
magnetic field:
\begin{equation}
  \vect{v}_{\mathrm{g}}
  = \frac{\partial\omega}{\partial\vect{k}}
  = \vA\,\hat{\vect{z}},
  \label{eq:alfven-group}
\end{equation}
showing that Alfv\'en waves transport energy along field lines, exactly
as in the non-relativistic theory.

\subsection{Comparison with Chandrasekhar~\S39}

Chandrasekhar~\cite{Chandrasekhar1961}, \S39, derives the Alfv\'en wave
dispersion for an incompressible, non-relativistic conducting fluid
permeated by a uniform magnetic field $\vect{H}$.  His result (in our
notation) is
\begin{equation}
  \omega^{2} = \frac{(\vect{k}\cdot\vect{H})^{2}}{4\pi\rho}
  = k^{2}\,V_{\!A}^{2}\,\cos^{2}\!\theta,
  \label{eq:alfven-chandrasekhar}
\end{equation}
with $V_{\!A}^{2} = H^{2}/(4\pi\rho)$.  This is recovered from our
relativistic result \eqref{eq:alfven-dispersion} as follows.

\subsection{Non-relativistic limit}

In the limit $v/c\to 0$:
\begin{enumerate}
  \item $\varepsilon \to \rho c^{2}$ (rest-mass energy dominates),
  \item $p \ll \rho c^{2}$,
  \item $b^{2} = B^{2} \ll 4\pi\rho c^{2}$ (i.e.\ $\vA \ll c$).
\end{enumerate}
Then
\begin{equation}
  \vA^{2}
  = \frac{b^{2}\,c^{2}}{4\pi(\varepsilon+p)+b^{2}}
  \;\longrightarrow\;
  \frac{B^{2}\,c^{2}}{4\pi\rho c^{2}}
  = \frac{B^{2}}{4\pi\rho}
  = V_{\!A}^{2},
  \label{eq:vA-NR-limit}
\end{equation}
recovering the classical Alfv\'en speed.  The dispersion relation
\eqref{eq:alfven-dispersion} then reduces identically to
\eqref{eq:alfven-chandrasekhar}.

\subsection{Ultra-relativistic regime}

In the opposite extreme, when $b^{2}\gg 4\pi(\varepsilon+p)$ (a
magnetically dominated plasma), the Alfv\'en speed saturates at
\begin{equation}
  \vA^{2} \to c^{2}
  \qquad\text{(magnetically dominated limit)}.
  \label{eq:vA-ultra}
\end{equation}
This is relevant for pulsar magnetospheres, relativistic jets, and
gamma-ray burst outflows where the magnetisation parameter
$\sigma_{\mathrm{mag}} = b^{2}/[4\pi(\varepsilon+p)] \gg 1$.  Even in
this extreme, the Alfv\'en speed respects $\vA < c$, as required by
causality.

\subsection{Polarisation and eigenmodes}

The Alfv\'en wave eigenmode has the polarisation
\begin{equation}
  \delta\vect{v} = \pm\,\frac{\delta\vect{B}}{\sqrt{4\pi\enthalpy^{\ast}_{0}}},
  \label{eq:alfven-eigenmode}
\end{equation}
where the $\pm$ corresponds to forward/backward propagating waves.
The magnetic and velocity perturbations are aligned (equipartition), and
the total pressure perturbation $\delta p + \delta(b^{2}/2)/(4\pi) = 0$.
These properties are identical to the non-relativistic Alfv\'en mode; the
relativistic modification enters only through the inertia
$\enthalpy^{\ast}_{0}$, which includes both rest-mass and magnetic
contributions.


%======================================================================
\section{Summary}
\label{sec:rmhd-summary}
%======================================================================

The relativistic MHD framework presented in this chapter provides the
foundation for all subsequent stability analyses.  The key results are:
\begin{enumerate}
  \item The ideal MHD condition $F^{\mu\nu}u_{\nu}=0$ leads to a
        well-defined energy-momentum tensor \eqref{eq:Tmunu-MHD} and a
        covariant induction equation \eqref{eq:induction-covariant}.
  \item All characteristic speeds---Alfv\'en, fast and slow
        magnetosonic---are automatically bounded by $c$
        (eq.~\ref{eq:vf-causal-proof}), ensuring causal propagation.
  \item Resistive effects are captured by the relativistic Ohm's law
        \eqref{eq:rel-Ohm}, with magnetic diffusivity
        $\eta_{\mathrm{m}}=c^{2}/(4\pi\sigma)$.
  \item The linearised perturbation theory yields the same modal
        structure (Alfv\'en, fast, slow) as the non-relativistic theory,
        but with relativistically corrected speeds.
  \item All results reduce to the classical Chandrasekhar framework in
        the limit $v/c\to 0$.
\end{enumerate}

\bigskip
\noindent
The next chapters will apply this framework to specific instability
problems: thermal convection in magnetised relativistic fluids
(the relativistic B\'enard problem), the relativistic
magnetorotational instability, and relativistic magnetic buoyancy
instabilities.
